\begin{center}
{\Large\bfseries ALGEBRAIC GEOMETRY AND ANALYTIC GEOMETRY}

\medskip
by Jean-Pierre \textsc{Serre}

\medskip
\rule{2.5cm}{0.4pt}
\end{center}

\tableofcontents

\GAGASection{Introduction}{gaga:introduction}

Let $\mathrm{X}$ be a projective algebraic variety defined over the field of complex numbers. The study of $\mathrm{X}$ may be approached from two points of view: the \textit{algebraic} point of view, in which one is concerned with the local rings at the points of $\mathrm{X}$ and with rational or regular maps from $\mathrm{X}$ to other varieties, and the \textit{analytic} point of view (sometimes called ``transcendental''), in which the notion of a holomorphic function on $\mathrm{X}$ plays the principal role. It is well known that the second point of view has proved particularly fruitful when $\mathrm{X}$ is nonsingular, since this hypothesis makes available all the resources of the theory of K\"ahler manifolds (harmonic forms, currents, cobordism, etc.).

In many questions the two points of view lead to essentially equivalent results, although by very different methods. For example, holomorphic differential forms at every point of $\mathrm{X}$ are precisely the rational differential forms that are everywhere ``of the first kind'' (still assuming that $\mathrm{X}$ is nonsingular); \textsc{Chow}'s theorem, according to which every closed analytic subspace of $\mathrm{X}$ is an algebraic variety, is another example of the same kind.

The principal aim of this paper is to extend this equivalence to \textit{coherent sheaves}. More precisely, we prove that coherent algebraic sheaves and coherent analytic sheaves correspond bijectively, and that the correspondence between these two categories of sheaves preserves cohomology groups (see \GAGAInternalRef{gaga:no:12}{no.~12} for the statements). We give various applications of these results, notably to the comparison between algebraic and analytic fibre spaces.

The first two sections are preliminary. In \GAGAInternalRef{gaga:section:1}{\S 1} we recall the definition and principal properties of ``analytic spaces.'' The definition we adopt is the one proposed by H. \textsc{Cartan} in \GAGACite{3}, except that H. \textsc{Cartan} restricted himself to \textit{normal} varieties, a restriction unnecessary for our purposes; a very similar definition was used by W.-L. \textsc{Chow} in his as yet unpublished work on this subject. In \GAGAInternalRef{gaga:section:2}{\S 2} we show how every algebraic variety $\mathrm{X}$ may be endowed with a structure of analytic space, and we give various elementary properties of that structure. Perhaps the most important is the fact that, if $\mathcal{O}_x$ (respectively $\mathcal{H}_x$) denotes the local ring (respectively the ring of germs of holomorphic functions) of $\mathrm{X}$ at $x$, then $\mathcal{O}_x$ and $\mathcal{H}_x$ have the same completion and consequently form a ``flat pair'' in the sense of the \GAGAInternalRef{gaga:appendix}{Appendix}, \GAGAInternalRef{gaga:def:4}{Definition~4}.

\GAGAInternalRef{gaga:section:3}{Section~3} contains the proofs of the theorems on coherent sheaves just mentioned. These proofs rest principally on the theory of coherent algebraic sheaves developed in \GAGACite{18} and on theorems A and B of \GAGACite{3}, Expos\'es XVIII--XIX; for completeness we reproduce the proofs of those theorems.

\GAGAInternalRef{gaga:section:4}{Section~4} is devoted to applications\footnote{We have left aside the applications to automorphic functions, for which we refer to \GAGACite{3}, Expos\'e XX.}: invariance of the \textsc{Betti} numbers under automorphisms of the field of complex numbers, \textsc{Chow}'s theorem, and comparison of algebraic and analytic fibre spaces having a given algebraic group as structure group. Our results on the last question are, moreover, quite incomplete: among all semisimple groups we can treat only the unimodular linear group and the symplectic group.

Finally, we require a number of results on local rings that do not appear explicitly in the literature; we have collected them in an \GAGAInternalRef{gaga:appendix}{Appendix}.
